% registry: LIUMAHADEVAN
\methodhead{LIUMAHADEVAN}{Liu and Mahadevan (2005): characteristic-plane criterion and life prediction}\label{sheet:LIUMAHADEVAN}

\block{Source} Y.~Liu, S.~Mahadevan, \emph{Int.\ J.\ Fatigue} 27 (2005) 790--800 \cite{liu2005},
Eqs.~(6), (19)--(25); Table~1; Figs.~2--5.

\block{Idea and scope} The fatigue fracture plane is the plane of maximum normal stress amplitude. The
critical plane lies at a material angle $\alpha$ from it, chosen so that the damage parameter meets
both fully reversed limits while the hydrostatic term is as small as possible; the critical plane
therefore depends on the material, not only on the stress state. The damage combines the normal,
shear and hydrostatic stress amplitudes; a normal mean on the critical plane multiplies the normal
amplitude by $1+\eta\,\sigma_{m,c}/\fm$, and a mean shear has no effect. Proposed for high-cycle
fatigue limits and lives of metals from ductile to extremely brittle ($\tm/\fm$ from 0.57 to 1.28),
under in-phase and out-of-phase loading with mean stresses.

\block{Inputs}
\begin{center}\small
\begin{tabular}{@{}llp{100mm}@{}}
\toprule
Symbol & Unit & Meaning\\\midrule
$\sigma_{a,c}$, $\sigma_{m,c}$ & MPa & normal stress amplitude and mean on the critical plane\\
$\tau_{a,c}$ & MPa & shear stress amplitude on the critical plane\\
$\sigma^H_a$ & MPa & hydrostatic stress amplitude, $I_{1,a}/3$ (Eq.~\eqref{eq:invariants})\\
$\fm$, $\tm$ & MPa & fully reversed bending (axial) and torsion fatigue limits\\
$s$ & -- & limit ratio, $s=\tm/\fm$ (or $t_N/f_N$ at finite life)\\
$\alpha$ & rad & angle between the fracture plane and the critical plane, Eq.~\eqref{eq:lm-alpha}\\
$\beta$, $k$ & --, -- & material constants of Eq.~\eqref{eq:lm-alpha} and Table~1\\
$\eta$, $\eta_N$ & -- & mean normal stress sensitivity, Eq.~\eqref{eq:lm-eta}, at the limit and at life $N_f$\\
$\sigma_f'$, $b$ & MPa, -- & axial Basquin coefficient and exponent, $f_N=\sigma_f'(2N)^b$ (workbook \texttt{sigma\_f}, \texttt{b\_fat})\\
$\tau_f'$, $b_0$ & MPa, -- & torsional Basquin coefficient and exponent, $t_N=\tau_f'(2N)^{b_0}$ (workbook \texttt{tau\_f}, \texttt{b0})\\
$f_N$, $t_N$ & MPa & axial and torsional fatigue strengths at life $N_f$\\
$N_f$ & cycles & fatigue life\\
$I$ & -- & error index, Eq.~(25)\\
\bottomrule
\end{tabular}
\end{center}

\block{Formulation} The criterion (Eq.~20) is
\begin{equation}
\sqrt{\left[\frac{\sigma_{a,c}\bigl(1+\eta\,\sigma_{m,c}/\fm\bigr)}{\fm}\right]^2
+\Bigl(\frac{\tau_{a,c}}{\tm}\Bigr)^2 + k\Bigl(\frac{\sigma^H_a}{\fm}\Bigr)^2} = \beta ,
\label{eq:lm}
\end{equation}
with, for $s=\tm/\fm$ (Eq.~19 and Table~1),
\begin{equation}
\begin{aligned}
s\le1:&\quad \cos2\alpha=\frac{-2+\sqrt{4-4(1/s^2-3)(5-1/s^2-4s^2)}}{2(5-1/s^2-4s^2)},\quad
\beta=\sqrt{\cos^2(2\alpha)\,s^2+\sin^2(2\alpha)},\quad k=0;\\
s>1:&\quad \alpha=0,\quad \beta=s,\quad k=9(s^2-1),
\end{aligned}
\label{eq:lm-alpha}
\end{equation}
and (Eq.~21)
\begin{equation}
\eta = \frac34 + \frac14\,\frac{\sqrt3-\fm/\tm}{\sqrt3-1}\quad(s\le1),\qquad \eta=1\quad(s>1).
\label{eq:lm-eta}
\end{equation}
The error index is $I=(\mathrm{LHS}-\mathrm{RHS})/\mathrm{LHS}$ (Eq.~25). \emph{Critical plane}: at
angle $\alpha$ from the plane of maximum normal stress amplitude; in the repository, the worst plane
on the cone of planes at $\alpha$ around the fracture-plane normal (see Repository choices); fracture
planes of equal amplitude: the worst is kept. \emph{Shear amplitude}: $\tau_{a,c}$ is the \texttt{rss}
of the harmonic shear components (MCE), by maintainer decision. The model is case-level and its
amplitude is fixed: the global selection of \texttt{docs/shear-amplitude.md} leaves it unchanged.

\block{Constants and calibration} Paper: $\alpha$, $\beta$, $k$ in closed form from $\tm/\fm$; $\eta$
calibrated on uniaxial mean-stress tests or taken from Eq.~(21); the life uses the material's uniaxial
and torsional S--N curves. Repository: the same closed forms and Eq.~(21); the S--N curves are the
workbook Basquin constants, or fitted by the pipeline on its $R=-1$ tension and torsion failures. No
master curve is fitted.

\block{Life} Native (Eqs.~23--24): the limits are replaced by the S--N strengths at life $N_f$,
\begin{equation}
\frac1\beta\sqrt{\Bigl[\sigma_{a,c}\Bigl(1+\eta_N\frac{\sigma_{m,c}}{f_N}\Bigr)\Bigr]^2
+\Bigl(\frac{f_N}{t_N}\Bigr)^2\tau_{a,c}^2 + k\,(\sigma^H_a)^2} = f_N ,
\label{eq:lm-life}
\end{equation}
with $\alpha$, $\beta$, $k$, $\eta$ evaluated at $s=t_N/f_N$, solved for $N_f$ by trial and error;
any S--N form may be used. The repository solves Eq.~\eqref{eq:lm-life} by bisection in $\log N$ with
$f_N=\sigma_f'(2N)^b$, $t_N=\tau_f'(2N)^{b_0}$; the comparison uses this native life (no DP
comparison).

\block{Repository choices}
\begin{itemize}
\item Plane rule in 3D (maintainer decision): the paper rotates the plane in 2D; here the critical
  plane is the worst plane on the cone at angle $\alpha$ around the fracture-plane normal. This meets
  both reversed limits exactly for $s\ge0.5$; below, a tilted plane of the cone exceeds the torsion
  limit, so $s<0.5$ raises \texttt{ValueError}.
\item Shear amplitude (maintainer decision): \texttt{rss} of the harmonic shear components, chosen on
  Fig.~2 (it matches hard steel at $0^\circ$, $60^\circ$, $90^\circ$; the longest chord gives 4.6\% at
  $90^\circ$).
\item Eq.~(19) is evaluated in the rationalised form $\cos2\alpha=-C/(1+\sqrt{1-AC})$, $A=5-1/s^2-4s^2$,
  $C=1/s^2-3$, which removes the $0/0$ at $s=1/2$ and $s=1$.
\item $\eta$ from Eq.~(21); experimental $\eta$ values are not used. The factor
  $1+\eta\sigma_{m,c}/f$ is clipped at zero.
\item Life (maintainer decision): native Eq.~(23) with Basquin curves; the bisection is kept to lives
  where $t_N/f_N\ge0.5$; a solution outside raises \texttt{LoadingScopeError}, and no solution below
  $10^{12}$ cycles gives an infinite life.
\item The index $\sqrt{\cdots}/\beta$ of Eq.~(20), 1 on the limit, is reported as a diagnostic only. No
  \texttt{\_ext} variant.
\end{itemize}

\block{Worked example} Hard steel of Papadopoulos et al.\ (1997): $\fm=313.9$, $\tm=196.2\MPa$,
$s=0.625$, so $\alpha=39.2^\circ$, $\beta=0.9875$, $k=0$. Papadopoulos Table~1 row~1
($\sigma_a=138.1$, $\tau_a=167.1\MPa$, $\delta=0^\circ$): index 0.975, $I=-2.5\%$; row~4
($\sigma_a=150.2$, $\tau_a=181.7\MPa$, $\delta=90^\circ$): index 1.056, $I=5.3\%$. Native life with
illustrative Basquin curves ($\sigma_f'=1200\MPa$, $b=-0.08$, $\tau_f'=700\MPa$, $b_0=-0.06$) at
$\sigma_a=400$, $\tau_a=230\MPa$: $N_f=1.29\cdot10^4$ in phase and $3.63\cdot10^4$ at $90^\circ$.

\block{Verification} Eq.~(19)/Table~1 values, including the removable singular points; both reversed
limits on the $\alpha$-cone for $0.5\le s\le1.28$; Eq.~(23) returns the input life on both S--N curves.
Figs.~2--5 (mean absolute error per phase on the Papadopoulos et al.\ \cite{papadopoulos1997} Table~1
data; the figures' Matake and McDiarmid curves equal that table's phase means): reproduced for hard
steel at $0^\circ$, $60^\circ$, $90^\circ$ and 30NCD16 at $60^\circ$. Strict xfails (unresolved): hard
steel at $30^\circ$, 42CrMo4, 34Cr4 and 30NCD16; the balanced $90^\circ$ out-of-phase rows come out
about 10\% non-conservative here and small in the paper, and no plane-rule or amplitude variant tried
closes the gap. Not reproducible: Table~4 (fracture angles, needs Macha's loads) and Figs.~10--13
(lives, data sets not given).
